\section*{Introduction}
\phantomsection\addcontentsline{toc}{section}{Introduction}

The previous chapter showed that current DR grading systems share three architectural
limitations: they ignore the relations between patients, they ignore the global
structure of the vascular tree, and they are rarely trained or evaluated with the
referral decision in mind. Each limitation calls for a family of methods that is well
established on its own but rarely used in medical image grading: graph neural networks for
relations between patients, topological data analysis for the structure of the vessels,
and ordinal prediction with calibrated probabilities for safety-oriented decisions.

This chapter introduces these methods and the tools used to evaluate them. We first
recall the deep-learning models used for medical images, from convolutional networks to
vision transformers and foundation models. We then present learning on graphs and the
inductive GraphSAGE model. Next, we introduce
persistent homology, which turns the shape of the vessel skeleton into numbers. We present
ordinal regression and the CORAL method, followed by uncertainty, calibration and
selective prediction, and by explainability. Finally, we define the evaluation metrics and
the statistical tests used in the experimental chapters.

\section{Deep Learning for Medical Image Analysis}
\label{sec:dl}

\subsection{Convolutional Neural Networks}

A convolutional neural network (CNN) processes an image through a stack of layers that
apply small learned filters over the whole image, followed by non-linearities and pooling
operations~\cite{lecun2015deep,goodfellow2016deep}. Early layers detect edges and textures,
deeper layers combine them into parts and objects. Because the same filter is applied at
every position, a CNN has far fewer parameters than a fully connected network and is
naturally robust to small translations, which suits images of lesions that can appear
anywhere on the retina.

The history of CNNs for image recognition is largely a history of depth. AlexNet showed in
2012 that a deep network trained on GPUs could outperform hand-crafted
features on ImageNet~\cite{krizhevsky2012imagenet,deng2009imagenet}. VGG networks stacked
small $3\times3$ filters~\cite{simonyan2015vgg}; Inception modules combined filters of
several sizes~\cite{szegedy2016inception}; residual networks made very deep models
trainable by adding shortcut connections~\cite{he2016resnet},
\begin{equation}
x_{l+1} = x_l + \mathcal{F}(x_l;\,W_l),
\label{eq:resnet}
\end{equation}
so that each block of \cref{eq:resnet} only learns a correction $\mathcal{F}$ to its input. DenseNet connected
every layer to all previous ones~\cite{huang2017densenet}, and EfficientNet scaled depth,
width and resolution together with a single coefficient~\cite{tan2019efficientnet}.
EfficientNet-B3 is the convolutional baseline used throughout this thesis. In medical
imaging, the same architectures, sometimes combined with segmentation networks such as
U-Net~\cite{ronneberger2015unet}, now dominate most tasks~\cite{litjens2017survey}, with
results that match specialists in dermatology~\cite{esteva2017dermatologist} and
ophthalmology~\cite{gulshan2016jama}.

\subsection{Attention and Vision Transformers}

Convolutions only look at a small neighbourhood at a time; relations between distant
regions of the image are captured indirectly, after many layers. The transformer
architecture~\cite{vaswani2017attention} replaces convolution by \emph{self-attention}, in
which every element of a sequence is updated with a weighted sum of all the others:
\begin{equation}
\mathrm{Attention}(Q,K,V) = \mathrm{softmax}\!\left(\frac{QK^\top}{\sqrt{d_k}}\right)V,
\label{eq:attention}
\end{equation}
where, in \cref{eq:attention}, the queries $Q$, keys $K$ and values $V$ are linear projections of the input. The
vision transformer (ViT)~\cite{dosovitskiy2021vit} applies this mechanism to an image cut
into patches of $16\times16$ pixels, each treated as a token. Transformers need more data
than CNNs to train from scratch, but scale better, and hierarchical variants such as the
Swin transformer~\cite{liu2021swin} bring back some of the locality of convolutions.
Attention weights are also used, with caution, as a form of explanation, since they show
which parts of the input the model relied on.

\subsection{Transfer Learning and Foundation Models}

Medical datasets are small compared with natural-image datasets, and labels are costly.
Transfer learning addresses this problem by initialising a model with weights learned on
another task, usually ImageNet classification, and fine-tuning it on the medical
task~\cite{pan2010transfer}. This practice is almost universal, although the features
learned on natural images are not always the most relevant for medical
images~\cite{raghu2019transfusion}.

Self-supervised learning removes the need for labels during pre-training: a network learns
by solving a task built from the data itself, such as recognising two augmented views of
the same image~\cite{chen2020simclr} or reconstructing masked image patches~\cite{he2022mae}.
Models pre-trained in this way on very large and diverse data, and then adapted to many
downstream tasks, are called \emph{foundation models}~\cite{bommasani2021foundation}.
RETFound~\cite{zhou2023retfound} is a vision transformer pre-trained with masked
autoencoding on $1.6$ million unlabelled retinal images. The systems of this thesis use
ImageNet-pre-trained EfficientNet backbones; replacing them by a retinal foundation model
is a direction for future work.

\subsection{Training Deep Networks}

Deep networks are trained by minimising a loss function with stochastic gradient descent
or one of its adaptive variants, of which Adam~\cite{kingma2015adam} is the most widely
used. Regularisation techniques such as dropout~\cite{srivastava2014dropout}, batch
normalisation~\cite{ioffe2015batchnorm} and data augmentation (random rotations, flips,
colour changes) reduce overfitting. For graded medical images, two further issues matter:
class imbalance, because severe grades are rare, and the ordinal nature of the labels.
Both are addressed in this thesis through the loss function (\cref{sec:coral}).

\section{Learning on Graphs}
\label{sec:graphs}

\subsection{Graphs and Message Passing}

A graph $G=(V,E)$ is a set of nodes $V$ connected by edges $E$. In a population graph,
each node is a patient (or a fundus image) described by a feature vector $x_v$, and an edge
between two nodes means that the two patients are similar. A graph neural network (GNN)
computes a representation $h_v$ of each node by repeatedly combining its own features with
those of its neighbours $\mathcal{N}(v)$. After $L$ rounds of this \emph{message passing},
the representation of a node summarises its $L$-hop neighbourhood, and a classifier
applied to $h_v$ can use information from similar patients to grade the current one.

Three architectures dominate the literature. The graph convolutional network
(GCN)~\cite{kipf2017gcn} averages neighbour features with fixed weights given by the graph
structure. The graph attention network (GAT)~\cite{velickovic2018gat} learns a weight
$\alpha_{ij}$ for each edge, so that the most informative neighbours count more. GraphSAGE,
presented below, samples a fixed number of neighbours and learns how to aggregate them.

Graph neural networks were first defined as recurrent models that propagate information
until a fixed point is reached~\cite{scarselli2009gnn}. Spectral approaches then defined
convolution on graphs through the eigenvectors of the graph
Laplacian~\cite{bruna2014spectral}, and fast polynomial approximations made them
practical~\cite{defferrard2016cheb}. Most current models can be written in the
\emph{message-passing} framework~\cite{gilmer2017mpnn,hamilton2020grl}, illustrated in
\cref{fig:mp}: at each layer, every node receives messages from its neighbours, aggregates
them, and updates its own representation.

% Figure 2.1 --- one round of message passing on a population graph
\begin{figure}[!htbp]
\centering
\resizebox{0.92\linewidth}{!}{%
\begin{tikzpicture}[
  font=\footnotesize,
  pt/.style={circle, draw=axisR, fill=axisRf, minimum size=7mm, inner sep=0},
  tg/.style={circle, draw=axisO, fill=axisOf, very thick, minimum size=8mm, inner sep=0},
  ed/.style={black!45, thick},
  msg/.style={-{Latex[length=2mm]}, axisR, very thick}
]
\begin{scope}
\node[tg] (v) at (0,0) {$v$};
\node[pt] (a) at (-1.6,1.1) {$u_1$};
\node[pt] (b) at (1.6,1.2) {$u_2$};
\node[pt] (c) at (-1.4,-1.3) {$u_3$};
\node[pt] (d) at (1.5,-1.2) {$u_4$};
\node[pt] (e) at (-3.0,0.1) {};
\node[pt] (g) at (3.0,0) {};
\foreach \x in {a,b,c,d} \draw[ed] (v) -- (\x);
\draw[ed] (a) -- (e); \draw[ed] (c) -- (e); \draw[ed] (b) -- (g); \draw[ed] (d) -- (g);
\node[below=12mm of v, align=center] {(a) patient $v$ and its neighbours\\in the population graph};
\end{scope}
\begin{scope}[xshift=7.2cm]
\node[tg] (v2) at (0,0) {$v$};
\node[pt] (a2) at (-1.6,1.1) {$u_1$};
\node[pt] (b2) at (1.6,1.2) {$u_2$};
\node[pt] (c2) at (-1.4,-1.3) {$u_3$};
\node[pt] (d2) at (1.5,-1.2) {$u_4$};
\foreach \x in {a2,b2,c2,d2} \draw[msg] (\x) -- (v2);
\node[draw=axisO, rounded corners=2pt, fill=white, align=center, font=\scriptsize] at (0,2.25)
  {$h_v^{(l)}=\sigma\bigl(W^{(l)}[\,h_v^{(l-1)}\,\Vert\,\mathrm{AGG}\{h_u^{(l-1)}\}\,]\bigr)$};
\node[below=12mm of v2, align=center] {(b) neighbours send their representations,\\which are aggregated and combined with $h_v$};
\end{scope}
\end{tikzpicture}}
\caption[One round of message passing on a population graph]{One round of message passing.
(a)~In a population graph, each node is a patient and edges connect similar patients.
(b)~Each neighbour sends its current representation; the messages are aggregated (for
example by a mean) and combined with the representation of $v$. After $L$ rounds, $h_v$
summarises the $L$-hop neighbourhood of the patient.}
\label{fig:mp}
\end{figure}


The choice of aggregation function determines what a GNN can distinguish: sum
aggregation is strictly more expressive than mean or max aggregation, as shown by the
graph isomorphism network (GIN)~\cite{xu2019gin}. Comprehensive surveys cover the many
variants proposed since~\cite{wu2021gnnsurvey,bronstein2017geometric}. In medicine, graphs
have been used to represent molecules, anatomical structures, brain connectivity and
populations of patients~\cite{ahmedt2021graphmed}. In a population graph, the edge
weights encode how similar two patients are, using image features or clinical
information~\cite{parisot2018disease}; the quality of this similarity is crucial, since a
graph that connects dissimilar patients propagates noise rather than signal.

\subsection{Inductive Inference with GraphSAGE}
\label{sec:sage}

GraphSAGE~\cite{hamilton2017graphsage} learns aggregation functions over sampled
neighbour sets:
\begin{align}
h^{(l)}_{\mathcal{N}(v)} &= \operatorname{MEAN}\bigl(\{\,h^{(l-1)}_u : u\in\mathcal{N}(v)\,\}\bigr),
\label{eq:sage-agg}\\
h^{(l)}_v &= \sigma\bigl(W^{(l)}\cdot\bigl[\,h^{(l-1)}_v \,\Vert\, h^{(l)}_{\mathcal{N}(v)}\,\bigr]\bigr).
\label{eq:sage-upd}
\end{align}
With $S=25$ samples per layer and $L=2$ layers, each prediction aggregates $S^L=625$
neighbour embeddings at $O(S^L)$ cost, independent of the size of the cohort. In a
$k$-nearest-neighbour graph with small $k$ (for example $k=8$ in DGTS, \cref{ch:dots}), a
node usually has fewer than $S$ neighbours, in which case all of them are used; $625$ is
therefore an upper bound. This inductiveness is a prerequisite for a screening service: a
new patient is graded by retrieving its neighbours from the reference index, without
retraining the graph (\cref{tab:latency}). GCN
in its original formulation~\cite{kipf2017gcn} is transductive and requires the full
graph at training time; GAT~\cite{velickovic2018gat} can be applied inductively (it was
evaluated on an inductive benchmark in the original paper), but it attends over complete
neighbourhoods, so its per-node cost grows with node degree (\cref{sec:gnnvariants}). For a screening service, where the population grows continuously, this difference is decisive: an inductive model can be deployed once and used for years, whereas a transductive one would have to be retrained whenever new patients join the graph.

\subsection{Choosing a Graph Architecture for Screening}
\label{sec:gnnvariants}

Three families of graph neural network were compared for clinical suitability.
\Cref{tab:gnn} compares their inference mode, per-patient cost, retraining need and
support for the per-prediction transparency required by Article~13 of the EU AI
Act~\cite{euaiact2024}.

\begin{table}[!htbp]
\centering
\caption[GNN variant comparison]{GNN variant comparison; inductive models with sampled
neighbourhoods suit a growing screening cohort. ``Art.~13'': support for per-prediction
transparency.}
\label{tab:gnn}
\small
\begin{tabular}{@{}lllP{3.3cm}c@{}}
\toprule
Model & Inference & Per-patient cost & Retraining for a new patient & Art.~13 \\
\midrule
GCN~\cite{kipf2017gcn} & Transductive & $O(|E|)$ per layer & Yes, full graph & Low \\
GAT~\cite{velickovic2018gat} & Inductive & $O(|\mathcal{N}(v)|)$ per layer & No & Medium \\
GraphSAGE~\cite{hamilton2017graphsage} & Inductive & $O(S^L)$, constant & No & Medium$^{\ast}$ \\
\bottomrule
\end{tabular}
\end{table}

Inductiveness is therefore a prerequisite for screening, which is why the graph-based
systems of this thesis use GraphSAGE. ($^{\ast}$) Although GraphSAGE has no attention
weights, the neighbours it aggregates are themselves an explanatory signal: DGTS returns
the retrieved similar cases with their verified grades (\cref{sec:audit}).

In practice, building a population graph requires three choices: the feature vector that
describes each patient, the similarity used to compare two patients, and the number $k$
of neighbours kept for each node. The first two decide which patients are considered
similar, and therefore what information the graph can propagate; the third controls a
trade-off between noise (too many weak neighbours) and isolation (too few). For cohorts of
tens of thousands of images, the $k$ nearest neighbours are found with approximate search
libraries such as FAISS~\cite{johnson2021faiss}, which keeps graph construction fast
enough for a screening service.

\section{Topological Data Analysis}
\label{sec:ph}

\subsection{From Vessel Skeletons to Point Clouds}

Topological data analysis (TDA) studies the \emph{shape} of data: how many connected
pieces it has, how many loops, and at which scale these features appear and disappear.
For retinal images, the object of interest is the vascular tree. After segmentation and
thinning, the vessels are reduced to a one-pixel skeleton, and the branch points of this
skeleton form a point cloud in the plane. Healthy vessels form a regular, tree-like
network; as DR progresses, capillaries close, abnormal shunts appear and new fragile vessels
grow, which changes the number of components and loops of the network.

\subsection{Simplicial Complexes and Homology}

Topology studies the properties of shapes that do not change under continuous
deformation, such as the number of connected pieces or of holes. To apply it to data, a
finite set of points must first be turned into a shape. A \emph{simplicial complex} does
this by gluing together simple building blocks: points (0-simplices), segments
(1-simplices), triangles (2-simplices) and their higher-dimensional
analogues~\cite{edelsbrunner2010computational}. The \emph{homology} of the complex then
counts its holes by dimension: the rank of $H_0$ is the number of connected components,
the rank of $H_1$ the number of independent loops (the first Betti number $\beta_1$), and
so on. For a vessel skeleton, $H_0$ reflects how fragmented the network is and $H_1$ how
many closed loops it contains.

\subsection{Persistent Homology}

For the branch points $P=\{p_1,\ldots,p_n\}$ of a vessel skeleton, the Vietoris--Rips
filtration adds a simplex $\sigma$ as soon as all pairwise distances within $\sigma$ are
at most $\varepsilon$. As $\varepsilon$ grows from $0$ to $\infty$, topological
features -- $H_0$ (connected components) and $H_1$ (independent loops) -- are born at scale
$b_j$ and die at $d_j$, with persistence $p_j=d_j-b_j$~\cite{edelsbrunner2002persistence,
carlsson2009topology}. The resulting persistence diagram $D$ summarises global structure
while being robust to local deformation: the stability
theorem~\cite{cohensteiner2007stability} bounds the change in the diagram by the change
in the filtering function,
\begin{equation}
d_B(D_f, D_g) \;\leq\; \lVert f-g\rVert_\infty ,
\label{eq:stability}
\end{equation}
where $d_B$ is the bottleneck distance. TOPORET compares diagrams with the closely related
Wasserstein distance $W_2$ (\cref{eq:dualsim}). This robustness makes the topological term of the
similarity less sensitive to small segmentation errors than a pixel-level comparison.

Persistence is what makes the method robust. A single threshold $\varepsilon$ would give
counts of components and loops that depend strongly on its value and on small errors of
segmentation; by following all thresholds at once, persistent homology keeps only the
features that survive over a wide range of scales. For a vessel skeleton, a short-lived
loop is typically an artefact of thinning, whereas a long-lived loop corresponds to a real
closed vascular circuit, such as an abnormal shunt or a neovascular frond.
\Cref{fig:filtration} illustrates the construction on a toy example. Persistent homology
was introduced as a way to separate robust topological features from
noise~\cite{edelsbrunner2002persistence,zomorodian2005computing}, and it is now the main
tool of topological data analysis~\cite{carlsson2009topology,chazal2021intro}.

% Figure 2.2 --- Vietoris--Rips filtration and persistence barcode on a toy point cloud
\begin{figure}[!htb]
\centering
\resizebox{\linewidth}{!}{%
\begin{tikzpicture}[font=\footnotesize]
\def\pts{(0,0),(1.2,0.1),(1.6,1.1),(0.8,1.8),(-0.2,1.1),(3.0,0.9)}
\foreach \k/\r/\lab in {0/0.08/{$\varepsilon=0$},1/0.62/{$\varepsilon_1$},2/0.8/{$\varepsilon_2$},3/1.05/{$\varepsilon_3$}}{
 \begin{scope}[xshift=\k*4.3cm]
  \foreach \p in \pts { \fill[axisSf, opacity=0.55] \p circle (\r); }
  \ifnum\k>0
   \draw[axisS, thick] (0,0)--(1.2,0.1)--(1.6,1.1)--(0.8,1.8)--(-0.2,1.1)--cycle;
  \fi
  \ifnum\k>2
   \fill[axisS, opacity=0.35] (0,0)--(1.2,0.1)--(1.6,1.1)--(0.8,1.8)--(-0.2,1.1)--cycle;
   \draw[axisS, thick] (1.6,1.1)--(3.0,0.9);
  \fi
  \foreach \p in \pts { \fill[black] \p circle (1.6pt); }
  \node at (1.4,-1.3) {\lab};
 \end{scope}}
% barcode
\begin{scope}[xshift=0cm, yshift=-4.6cm]
 \draw[-{Latex[length=2mm]}] (0,0) -- (16.5,0) node[right] {scale $\varepsilon$};
 \foreach \x/\l in {0/0, 4.3/{\varepsilon_1}, 8.6/{\varepsilon_2}, 12.9/{\varepsilon_3}} \draw (\x,0.08)--(\x,-0.08) node[below] {$\l$};
 \node[left] at (-0.2,1.55) {$H_0$};
 \foreach \y/\e in {1.1/4.3,1.3/4.3,1.5/4.3,1.7/4.3} \draw[axisR, line width=2pt] (0,\y) -- (\e,\y);
 \draw[axisR, line width=2pt] (0,1.9) -- (12.9,1.9);
 \draw[axisR, line width=2pt, -{Latex[length=2mm]}] (0,2.1) -- (16.3,2.1);
 \node[left] at (-0.2,0.55) {$H_1$};
 \draw[failred, line width=2.5pt] (4.3,0.55) -- (11.5,0.55);
 \node[right, font=\scriptsize, text=failred] at (11.6,0.55) {one persistent loop};
\end{scope}
\end{tikzpicture}}
\caption[Vietoris--Rips filtration and persistence barcode]{Vietoris--Rips filtration on a
toy point cloud of six points and the resulting barcode. As the scale $\varepsilon$ grows,
balls around the points merge: components ($H_0$ bars) die when they join, a loop ($H_1$
bar, red) is born when the five points on the left close a cycle and dies when the cycle is
filled. Long bars correspond to robust structure, short bars to noise. In this thesis the
points are the branch points of the vessel skeleton.}
\label{fig:filtration}
\end{figure}


\subsection{From Diagrams to Features}

A persistence diagram is a set of points of variable size, which cannot be fed directly to
a standard learning model. Several vector representations have been proposed: persistence
landscapes~\cite{bubenik2015landscapes}, persistence images~\cite{adams2017images}, learned
layers that process diagrams inside a neural
network~\cite{hofer2017topological,carriere2020perslay}, and simple summary statistics such
as total persistence, number of features or the longest bar. Efficient software makes these
computations practical on large datasets~\cite{otter2017roadmap,maria2014gudhi,
tauzin2021giotto}. This thesis uses a small vector of interpretable summary statistics
(\cref{eq:tda}), which can be related directly to vascular pathology.

\section{Ordinal Prediction}
\label{sec:coral}

\subsection{Why DR Grading Is an Ordinal Problem}

The five ICDR grades are ordered: Grade~3 is worse than Grade~2, which is worse than
Grade~1. A standard classifier ignores this order and treats a confusion between Grade~0
and Grade~4 exactly like a confusion between Grade~3 and Grade~4, although the first is a
serious clinical error and the second is a disagreement that also occurs between experts.
Ordinal regression methods decompose the problem into $K-1$ binary questions of the form
``is the grade at least $k$?'' and combine their answers.

\subsection{Families of Ordinal Methods}

Ordinal regression has a long history in statistics and machine
learning~\cite{gutierrez2016ordinal}. A simple and general approach transforms a
$K$-class ordinal problem into $K-1$ binary problems, ``is the grade greater than $k$?'',
and combines their probabilities~\cite{frank2001ordinal}. Deep versions of this idea train
a network with $K-1$ binary outputs~\cite{niu2016ordinal}, but nothing guarantees that the
estimated probabilities decrease with $k$, so the model can produce inconsistent rankings.
Other works keep a standard classifier and change the loss, for example by directly
optimising the weighted kappa~\cite{delatorre2018kappa}. CORAL and CORN, presented next,
are the two most widely used deep ordinal methods with a rank-consistency property.

\subsection{CORAL and CORN}

CORAL~\cite{cao2020coral} obtains rank consistency from the architecture rather than from
a loss penalty. All $K-1$ binary tasks share the same weight vector $w\in\mathbb{R}^d$ and
differ only by their biases $b_1,\ldots,b_{K-1}$, so that
$\hat P(Y\geq k)=\sigma(w^{\top}h+b_k)$. Cao et al.\ show that the biases minimising the
CORAL loss are ordered, $b_1\geq\cdots\geq b_{K-1}$; the ordering can also be enforced
explicitly during training by the projection
\begin{equation}
b_k \leftarrow \min(b_k,\, b_{k-1} - \epsilon), \qquad k=2,\ldots,K-1 .
\label{eq:coral-proj}
\end{equation}
With ordered biases, $\hat P(Y\geq 1)\geq \hat P(Y\geq 2)\geq\cdots$ holds for every input:
the cumulative probabilities cannot cross. CORN~\cite{shi2023corn} obtains rank
consistency differently: its $K-1$ binary tasks estimate conditional probabilities
$\hat P(Y>k\mid Y>k-1)$, and the cumulative probabilities are their chain-rule product,
which cannot increase with $k$ provided the outputs are decoded in this way.

\Cref{fig:coral} summarises the resulting architecture. The single shared weight vector
is what makes the guarantee possible: all binary questions use the same score $s_i$, and
only the thresholds $b_k$ move along it.

% CORAL head schematic
\begin{figure}[!tb]
\centering
\resizebox{0.95\linewidth}{!}{%
\begin{tikzpicture}[font=\footnotesize,
  box/.style={draw, rounded corners=2pt, align=center, inner sep=3pt, minimum height=0.75cm},
  arr/.style={-{Latex[length=2mm]}, thick}]
\node[box, fill=black!6, draw=black!55, minimum height=3.8cm, minimum width=1.7cm] (z) at (0,0) {Embedding\\$z_i$};
\node[box, fill=axisRf, draw=axisR, minimum height=3.8cm, minimum width=2.2cm] (w) at (3.2,0) {Shared\\weights $w$\\[4pt]$s_i = w^\top z_i$};
\node[font=\bfseries\scriptsize] at (6.9,2.35) {Logits};
\node[font=\bfseries\scriptsize] at (10.3,2.35) {$\hat P(Y\geq k)$};
\foreach \k/\y in {1/1.35,2/0.45,3/-0.45,4/-1.35}{
  \node[box, fill=axisOf, draw=axisO, minimum width=2.4cm] (g\k) at (6.9,\y) {$g_i^{(\k)} = s_i + b_\k$};
  \node[box, fill=white, draw=black!55, minimum width=2.2cm] (p\k) at (10.3,\y) {$\sigma\big(g_i^{(\k)}\big)$};
  \draw[arr] (w.east) -- (g\k.west);
  \draw[arr] (g\k) -- (p\k);
}
\draw[arr] (z) -- (w);
\node[box, fill=okf, draw=okgreen, minimum height=3.8cm, minimum width=3.2cm] (y) at (14.2,0) {Grade $\hat y_i$\\[4pt]number of $k$\\with $\hat P(Y\geq k)>0.5$};
\foreach \k in {1,2,3,4} \draw[arr] (p\k.east) -- (y.west |- p\k.east);
\node[font=\scriptsize, text=axisO, align=center] at (6.9,-2.55) {$b_1\geq b_2\geq b_3\geq b_4$\\enforced after each update};
\node[font=\scriptsize, align=center] at (10.3,-2.55) {decreasing in $k$\\$\Rightarrow 0$ rank violations};
\end{tikzpicture}}
\caption[The CORAL ordinal head]{The CORAL ordinal head. A single weight vector $w$ is shared
by the $K-1=4$ binary questions ``is the grade at least $k$?''; only the biases $b_k$ differ.
Because the biases are kept ordered, the estimated probabilities always decrease with $k$,
whatever the input and the weights, and the grade is simply the number of questions
answered ``yes''.}
\label{fig:coral}
\end{figure}


Rank consistency matters beyond elegance. A model that gives a higher probability to
``at least Grade~4'' than to ``at least Grade~3'' produces an output that has no clinical
meaning and cannot be explained to a grader. It also makes the choice of a referral rule
ambiguous, since different thresholds on the cumulative probabilities can lead to
contradictory decisions. With CORAL, the cumulative probabilities form a proper survival
function over the grades, and the grade is decoded by counting the binary tasks whose
probability exceeds $0.5$, $\hat y=\sum_k\mathbb{1}[\hat P(Y\geq k)>0.5]$; this is the
decoding used by DOTS (\cref{ch:dots}).

\section{Uncertainty, Calibration and Selective Prediction}
\label{sec:calib}

\subsection{Calibration}

The Expected Calibration Error (ECE)~\cite{naeini2015ece,guo2017calibration} partitions the $n$
predictions into $M$ confidence bins $B_m$ and measures the gap between accuracy and
confidence:
\begin{equation}
\mathrm{ECE} = \sum_{m=1}^{M} \frac{|B_m|}{n}\,\bigl|\operatorname{acc}(B_m) - \operatorname{conf}(B_m)\bigr| .
\label{eq:ece}
\end{equation}
An ECE of $0.031$, the value reached by DOTS after temperature scaling (\cref{ch:dots}),
means that on average its predicted confidence departs from empirical accuracy by about
$3$ percentage points. Temperature scaling~\cite{guo2017calibration} divides the logits
by a single scalar $T$ fitted on validation data, which changes the confidence but not the
predicted grade.

\subsection{Sources of Uncertainty}

The uncertainty of a prediction has two sources~\cite{kendall2017uncertainties}.
\emph{Aleatoric} uncertainty comes from the data itself: an image may be blurred, or the
lesions may be genuinely at the boundary between two grades, and more training data will
not remove this ambiguity. \emph{Epistemic} uncertainty comes from the model: it is high
for inputs unlike those seen during training, and decreases with more data. Estimating
epistemic uncertainty usually requires several predictions per input, obtained with Monte
Carlo dropout~\cite{gal2016dropout} or with an ensemble of independently trained
networks~\cite{lakshminarayanan2017ensembles}. In medical decision making, reporting
uncertainty is increasingly seen as a requirement rather than an
option~\cite{begoli2019uncertainty}. The systems of this thesis report calibrated
single-pass probabilities; ensemble or Monte Carlo estimates of epistemic uncertainty are
not used.

\subsection{Selective Prediction and Abstention}

The idea of letting a classifier reject uncertain inputs is old: Chow showed in 1970
that, for a given cost of errors and of rejections, the optimal rule rejects the inputs
whose highest class probability is below a threshold~\cite{chow1970reject}. A calibrated model knows, to some extent, when it is likely to be wrong. Selective
prediction uses this information to let the model \emph{abstain}: when the uncertainty of a
prediction exceeds a threshold, the case is deferred to a human expert instead of being
graded automatically~\cite{geifman2017selective}. Two quantities then describe the system:
the \emph{coverage}, i.e.\ the fraction of cases graded automatically, and the error rate on
the cases that are graded. Lowering the threshold defers more cases and reduces errors, at
the cost of more work for human graders. Selective prediction is not implemented in the
systems of this thesis; building it on the calibrated probabilities of DOTS is one of the
perspectives of the General Conclusion.

\section{Explainability}
\label{sec:xai}

A clinician who receives an automated grade needs to know why it was given, and European
regulation requires that high-risk AI systems be transparent to their users. Explanation
methods for image classifiers fall into two groups. \emph{Post-hoc} methods explain a
trained model from outside: class activation maps~\cite{zhou2016cam} and
Grad-CAM~\cite{selvaraju2017gradcam} highlight the image regions that most influenced the
decision, integrated gradients~\cite{sundararajan2017ig} attribute the prediction to input
pixels, and model-agnostic methods such as LIME~\cite{ribeiro2016lime} and
SHAP~\cite{lundberg2017shap} approximate the model locally with a simple one. For DR,
displaying such explanations has been shown to help graders~\cite{sayres2019integrated}.
Their limitation is that they approximate the model and may not reflect its actual
reasoning. \emph{Intrinsic} explanations, by contrast, come from quantities that the model
computes to make its decision, such as attention weights or interpretable features.
Intrinsic signals are not automatically faithful either: attention weights, in
particular, are not guaranteed to measure the causal importance of each input. The
systems of this thesis rely on such intrinsic explanatory signals -- the similar
patients retrieved by the graph with their verified grades, the topological features of
the vessel network and calibrated probabilities (\cref{sec:audit}) -- whose faithfulness was
not evaluated in this thesis.

\section{Evaluation Metrics and Statistical Tests}
\label{sec:metrics}

\subsection{Metrics}

The following metrics are used throughout the experimental chapters.
\begin{description}[style=nextline,leftmargin=1.2em]
\item[Quadratic weighted kappa (QWK)] Agreement between predicted and reference grades,
corrected for chance and weighted by the squared distance between grades:
\begin{equation}
\kappa_w = 1-\frac{\sum_{i,j} w_{ij}\,O_{ij}}{\sum_{i,j} w_{ij}\,E_{ij}},\qquad
w_{ij}=\frac{(i-j)^2}{(K-1)^2},
\label{eq:qwk}
\end{equation}
where, in \cref{eq:qwk}, $O$ is the observed confusion matrix and $E$ the matrix expected by
chance~\cite{cohen1968weighted}. QWK is the standard accuracy metric of DR grading
challenges.
\item[Area under the ROC curve (AUC)] Ability to separate referable from non-referable
images, independently of the decision threshold~\cite{hanley1982auc}.
\item[Expected calibration error (ECE)] Gap between confidence and accuracy
(\cref{eq:ece}).
\item[Severe-DR non-referral rate (Sev-NR)] Fraction of Grade~3 images predicted as No DR
or Mild, grades that do not trigger referral (\cref{eq:nr}). It is the primary safety
metric of this thesis, with a threshold of $2\%$.
\item[PDR non-referral rate (PDR-NR)] The same quantity for Grade~4.
\item[Rank violations] Fraction of predictions for which the cumulative probabilities
$\hat P(Y\geq k)$ are not decreasing in $k$.
\end{description}

\subsection{Statistical Tests}

Each hypothesis of the thesis is tested with an explicit statistical procedure. A paired
$t$-test compares two systems evaluated on the same cross-validation folds. A one-way
analysis of variance (ANOVA), followed by Tukey's post-hoc test~\cite{tukey1949}, tests
whether a feature differs across the five grades. Wilcoxon and McNemar tests compare paired
folds or paired predictions. Bootstrap resampling~\cite{efron1993bootstrap} gives
confidence intervals for QWK, and exact Clopper--Pearson intervals are used for the
non-referral rates, which are proportions of small counts. Effect sizes are reported with Cohen's $d$ and $\eta^2$~\cite{cohen1988}. The
formal statements are collected in \cref{app:stats}.

\section{Summary of the Methods}
\label{sec:methods-summary}

\Cref{tab:methods} summarises where each method introduced in this chapter is used in the
rest of the thesis and which limitation it addresses.

Two observations follow from \cref{tab:methods}. First, each limitation is addressed by a
method with a formal property rather than by a heuristic: inductiveness for the graph,
stability for persistence diagrams, rank consistency for CORAL and calibration for the
confidence scores. Second, the methods are complementary rather than redundant: the graph
and the topology act on the representation of a patient, whereas the ordinal head and
calibration act on the decision taken from that representation. The experimental chapters
test whether this complementarity holds in practice.

\begin{table}[!htbp]
\centering
\caption{Methods introduced in this chapter and their role in the thesis.}
\label{tab:methods}
\small
\begin{tabular}{@{}P{4.2cm} P{6.2cm} c c@{}}
\toprule
Method & Role in the thesis & Limitation & Chapters \\
\midrule
EfficientNet-B3/B4 & Visual features of the fundus image & -- & 3--5 \\
GraphSAGE & Inductive population graph; grading of a new patient without retraining & L1 & 3--5 \\
Persistent homology & Descriptors of the vessel skeleton (TOPORET) and of feature-space neighbourhoods (DGTS) & L2 & 3--5 \\
CORAL & Rank-consistent ordinal head of DOTS & L3 & 5 \\
Temperature scaling & Calibrated confidence scores & L3 & 5 \\
QWK, AUC, ECE, Sev-NR & Accuracy, calibration and safety evaluation & -- & 3--5 \\
\bottomrule
\end{tabular}
\end{table}

\section*{Conclusion}
\phantomsection\addcontentsline{toc}{section}{Conclusion}

This chapter introduced the tools of the thesis: convolutional features, the inductive
GraphSAGE model, persistent homology of the vessel skeleton, CORAL ordinal regression and
temperature scaling, together with metrics that separate accuracy, calibration and safety.
The next chapter tests these tools on simpler settings before applying them to the retina.
